\section{\new{Deployment, Interoperability, and Ethical Considerations}}

\subsection{\new{Deployment considerations}}

\new{Realistic deployment of CLINES depends on factors outside the extraction algorithm itself. First, \emph{input}: CLINES assumes de-identified, plain-text notes, and producing that input from raw EHR sources requires PHI de-identification, format normalization, and section/document segmentation, which are institutionally heterogeneous and outside the present scope. Second, \emph{compute}: the open-weight Llama-3.1-405B backbone was served in FP8 precision on 8$\times$H100 GPUs---the only single-node configuration that accommodates a 405B model---which entails a quality trade-off relative to FP16 and an infrastructure footprint beyond the reach of many small and medium centers; the hosted-API backbones substantially lower this barrier (Figure~\ref{fig:robustness-cost}B). Third, \emph{governance}: the hosted-API backbones (GPT-4o, o3-mini) were run via HMS Information Technology's managed Azure AI service (HMS Azure AI), an HMS-administered Microsoft Azure OpenAI deployment scoped to Harvard data-security Level~4 projects and operated on de-identified text only (HIPAA-regulated data is excluded from this platform); institutions handling identifiable PHI should select a corresponding HIPAA-compliant cloud configuration, and institutions with EU data-residency obligations should prefer an in-region deployment or a locally hosted open-weight backbone.}

\subsection{\new{Interoperability with production terminologies}}

\new{CLINES emits UMLS CUIs paired with preferred terms. UMLS is an ontology-grounded intermediate representation rather than a production interchange format: downstream use in FHIR resources or common data models requires mapping CUIs onward to the appropriate target vocabularies---SNOMED~CT (via UMLS cross-source mappings), LOINC (laboratory observables), or RxNorm (medications)---and OMOP-CDM ingestion via its standardized vocabularies, with attribute-field standardization. We therefore describe CLINES output as an ontology-grounded intermediate representation suitable for downstream mapping into production terminologies, rather than as a directly FHIR-ready artifact, and we recommend retaining a SNOMED~CT intermediate layer in deployments that target FHIR.}

\subsection{\new{Ethical considerations in demographic-attribute extraction}}

\new{CLINES includes a prompt that can extract demographic attributes (e.g., race, ethnicity, sex, and geographic identifiers). These fields are included to support downstream cohort discovery and health-equity auditing, not model personalization. We have not evaluated whether extraction accuracy varies across demographic subgroups, and we explicitly recommend a subgroup-stratified evaluation as a required pre-deployment step for any institutional adopter. Downstream users bear responsibility for applying appropriate access controls when demographic fields are stored or analyzed.}

\subsection{\new{Extended limitations and future-work directions}}

\new{For each main-text limitation we note a path forward: prompt sensitivity can be addressed with automatic prompt-search and optimization workflows; residual hallucination can be reduced with retrieval-grounded extraction that cites source spans and thereby enables provenance checking; and annotation variability can be narrowed by extending double-annotation to a larger validation cohort. Three further limitations bound the present claims. (v) The evaluation covers critical-care, COVID-era, and oncology notes but not several common note types (discharge summaries, operative and radiology reports), so generalizability beyond the evaluated distributions is untested; the bootstrap and significance analyses were further limited to the three consensus-annotated corpora. (vi) The relation-extraction module, though deployed, is not benchmarked here because gold-standard relation annotations are unavailable. (vii) We did not exhaustively optimize the single-prompt and chain-of-thought baselines, in order to avoid an unfair comparison in either direction; their reported performance should therefore be read as a reasonable rather than a maximally-tuned reference.}

\new{Further directions include: hybrid per-stage model routing, in which a lightweight model performs entity extraction while a stronger model is reserved for reasoning-intensive steps (the modular design already supports per-stage backbone selection); a formal output-consistency evaluation across repeated runs, given the stochastic decoding used here; retrieval-grounded extraction with span-level provenance to mitigate hallucination; subgroup-stratified fairness evaluation of demographic extraction; benchmarking of the relation-extraction module once gold relation annotations become available; expansion to additional note types (discharge summaries, operative and radiology reports); and tighter integration with OMOP-CDM and FHIR via SNOMED~CT, RxNorm, and LOINC mapping.}
